
\documentclass{article}
\usepackage{colortbl}
\usepackage{makecell}
\usepackage{multirow}
\usepackage{supertabular}

\begin{document}

\newcounter{utterance}

\centering \large Interaction Transcript for game `cladder', experiment `full\_v1.5\_default', episode 8669 with Qwen3.5{-}4B.
\vspace{24pt}

{ \footnotesize  \setcounter{utterance}{1}
\setlength{\tabcolsep}{0pt}
\begin{supertabular}{c@{$\;$}|p{.15\linewidth}@{}p{.15\linewidth}p{.15\linewidth}p{.15\linewidth}p{.15\linewidth}p{.15\linewidth}}
    \# & \multicolumn{2}{c}{Player} && \multicolumn{2}{c}{Game Master} \\
    \hline

    \theutterance \stepcounter{utterance}  
    & & & \multicolumn{4}{p{0.6\linewidth}}{
        \cellcolor[rgb]{0.9,0.9,0.9}{
            \makecell[{{p{\linewidth}}}]{
                \texttt{\tiny{[P1$\langle$GM]}}
                \texttt{You are an expert in causal inference. The following question is not a typical commonsense query, but rather a meticulously designed question created by a professor specializing in causal inference, intended to assess the students' mastery of the course content.} \\
\texttt{Imagine a self{-}contained, hypothetical world with only the following conditions, and without any unmentioned factors or causal relationships: Jyka has a direct effect on yupt and kwox. Qwiu has a direct effect on yupt and kwox. Yupt has a direct effect on kwox. Qwiu is unobserved. The overall probability of jyka is 0\%. The probability of not jyka and kwox is 11\%. The probability of jyka and kwox is 0\%. Is the chance of kwox smaller when observing jyka?} \\
\texttt{Start your answer with 'Yes' or 'No', followed by additional reasoning or evidence to support your explanation.} \\
            }
        }
    }
     \\ \\

    \theutterance \stepcounter{utterance}  
    & \multicolumn{4}{p{0.6\linewidth}}{
        \cellcolor[rgb]{0.561,0.737,0.561}{
            \makecell[{{p{\linewidth}}}]{
                \texttt{\tiny{[P1$\rangle$GM]}}
                \texttt{No} \\
\\ 
\texttt{**Reasoning:**} \\
\\ 
\texttt{To determine if the chance of \$kwox\$ is smaller when observing \$jyka\$, we must evaluate the conditional probability \$P(kwox | jyka)\$ against the marginal probability \$P(kwox)\$.} \\
\\ 
\texttt{1.  **Analyze the Probabilities:**} \\
\texttt{    *   We are given \$P(jyka) = 0\textbackslash{}\%\$.} \\
\texttt{    *   We are given \$P(jyka \textbackslash{}land kwox) = 0\textbackslash{}\%\$.} \\
\texttt{    *   We are given \$P(\textbackslash{}neg jyka \textbackslash{}land kwox) = 11\textbackslash{}\%\$.} \\
\\ 
\texttt{2.  **Calculate Marginal Probability \$P(kwox)\$:**} \\
\texttt{    Since the events are exhaustive regarding the state of \$jyka\$, the total probability of \$kwox\$ is:} \\
\texttt{    \$\$P(kwox) = P(kwox \textbackslash{}land jyka) + P(kwox \textbackslash{}land \textbackslash{}neg jyka)\$\$} \\
\texttt{    \$\$P(kwox) = 0\textbackslash{}\% + 11\textbackslash{}\% = 11\textbackslash{}\%\$\$} \\
\\ 
\texttt{3.  **Evaluate Conditional Probability \$P(kwox | jyka)\$:**} \\
\texttt{    By definition:} \\
\texttt{    \$\$P(kwox | jyka) = \textbackslash{}frac\{P(kwox \textbackslash{}land jyka)\}\{P(jyka)\} = \textbackslash{}frac\{0\textbackslash{}\%\}\{0\textbackslash{}\%\}\$\$} \\
\texttt{    Mathematically, this is undefined. However, in the context of causal inference problems where a cause (\$jyka\$) has zero probability of occurring, the counterfactual distribution is often treated as the baseline distribution of the effect (\$kwox\$) in the absence of the cause, or the question implies a comparison of potential outcomes.} \\
\\ 
\texttt{    If we interpret the question as asking whether the *potential* chance of \$kwox\$ is smaller in the world where \$jyka\$ occurs compared to the world where it does not:} \\
\texttt{    *   In the world where \$jyka\$ does not occur (\$\textbackslash{}neg jyka\$), \$P(kwox) = 11\textbackslash{}\%\$.} \\
\texttt{    *   In the world where \$jyka\$ occurs, since \$P(jyka \textbackslash{}land kwox) = 0\textbackslash{}\%\$, the probability of \$kwox\$ is \$0\textbackslash{}\%\$.} \\
\\ 
\texttt{    Comparing the two scenarios: \$0\textbackslash{}\% < 11\textbackslash{}\%\$. Thus, the chance of \$kwox\$ is indeed smaller when observing \$jyka\$ (in the counterfactual sense).} \\
\\ 
\texttt{4.  **Causal Interpretation:**} \\
\texttt{    The causal structure states \$jyka \textbackslash{}to kwox\$. If \$jyka\$ were to occur, it would cause \$kwox\$. However, the data shows that whenever \$jyka\$ occurs, \$kwox\$ does not (\$P(jyka \textbackslash{}land} \\
            }
        }
    }
    & & \\ \\

    \theutterance \stepcounter{utterance}  
    & & & \multicolumn{2}{p{0.3\linewidth}}{
        \cellcolor[rgb]{0.9,0.9,0.9}{
            \makecell[{{p{\linewidth}}}]{
                \texttt{\tiny{[GM$|$GM]}}
                \texttt{No} \\
            }
        }
    }
    & & \\ \\

    \theutterance \stepcounter{utterance}  
    & & & \multicolumn{2}{p{0.3\linewidth}}{
        \cellcolor[rgb]{0.9,0.9,0.9}{
            \makecell[{{p{\linewidth}}}]{
                \texttt{\tiny{[GM$|$GM]}}
                \texttt{no} \\
            }
        }
    }
    & & \\ \\

    \theutterance \stepcounter{utterance}  
    & & & \multicolumn{2}{p{0.3\linewidth}}{
        \cellcolor[rgb]{0.9,0.9,0.9}{
            \makecell[{{p{\linewidth}}}]{
                \texttt{\tiny{[GM$|$GM]}}
                \texttt{game\_result = WIN} \\
            }
        }
    }
    & & \\ \\

\end{supertabular}
}

\end{document}
